% !Mode:: "TeX:UTF-8"
\chapter{数据驱动范数最优迭代学习控制}

\section{预备知识与问题描述}

考虑一个最小实现（能控能观）的离散线性时不变系统：
\begin{equation}
	\begin{aligned}
		x_k(t+1) &= A x_k(t) + B u_k(t), \\
		y_k(t) &= C x_k(t),\\
		x_k(0) &= x_0, \ \forall k\in\mathbb{N}
	\end{aligned}
	\label{eq:sys}
\end{equation}
其中，$n \geq 0$，$m, p \geq 1$，$A \in \mathbb{R}^{n \times n}$，$B \in \mathbb{R}^{n \times m}$，$C \in \mathbb{R}^{p \times n}$，$D \in \mathbb{R}^{p \times m}$，\(u_k(t)\in\mathbb{R}^m\)，\(y_k(t)\in\mathbb{R}^p\)，\(x_k(t)\in\mathbb{R}^n\)。$u_k(t)$、$y_k(t)$ 和 $x_k(t)$ 分别为系统输入、输出和状态，系统的阶数为 \(n\)，输入维数为 \(m\)，输出维数为 \(p\)。$k \in \mathbb{N}$ 为迭代指标，$t \in [0, T]$ 为时间指标。系统 \eqref{eq:sys} 在有限时间区间内重复运行。当时间到达 $T$ 时，时间 $t$ 重置为 0，状态重置为相同的初始状态 $x_k(0) = x_0, \forall k \in \mathbb{N}$，然后开始下一次迭代。

我们采用“提升形式”表示（Hätönen et al., 2004），并且假设相对度为1，即 $\operatorname{rank}(CB) = p$，系统输入、输出以及期望轨迹定义为
\begin{align}
	\mathbf{u}_k &= [u_k(0), u_k(1),\dots,u_k(T-1)]^\top \in \mathbb{R}^{mT},\\
	\mathbf{y}_k &= [y_k(1), y_k(2),\dots,y_k(T)]^\top \in \mathbb{R}^{pT},\\
	\mathbf{r} &= [r(1), r(2),\dots,r(T)]^\top \in \mathbb{R}^{pT}
	\label{eq:lifted}
\end{align}
则系统输入输出关系可写为：
\begin{equation}
	\mathbf{y}_k = \mathbf{G} \mathbf{u}_k + \mathbf{d}
\end{equation}
其中，系统矩阵 $\mathbf{G}$ 为 Toeplitz 矩阵：
\begin{equation}
	\mathbf{G} = \begin{bmatrix}
		CB & 0 & \cdots & 0 \\
		CAB & CB & \cdots & 0 \\
		\vdots & \vdots & \ddots & \vdots \\
		CA^{T-1}B & CA^{T-2}B & \cdots & CB
	\end{bmatrix}
\end{equation}
而 $\mathbf{d}$ 为系统初始响应：
\begin{equation}
	\mathbf{d} = \big[ CAx_0, CA^{2}x_0, \dots, CA^{T}x_0 \big]^\top
\end{equation}
定义相应的跟踪误差 $\mathbf{e}_k = \mathbf{r} - \mathbf{y}_k$，则
\begin{equation}
	\mathbf{e}_k = \mathbf{r} - \mathbf{y}_k = \mathbf{r} - \mathbf{G} \mathbf{u}_k - \mathbf{d}
	\label{eq:tracking_error}
\end{equation} 
输入 $\mathbf{u}_k$ 和输出 $\mathbf{y}_k$ 分别定义在 Hilbert 空间 $\mathbf{U} = \mathbb{R}^{mT}$ 和 $\mathbf{Y} = \mathbb{R}^{pT}$ 上，其内积及诱导范数分别为
\begin{equation}
	\langle \mathbf{u}_k, \tilde{\mathbf{u}_k} \rangle_R = \mathbf{u}_k^\top R \tilde{\mathbf{u}_k}, \quad \langle \mathbf{y}_k, \tilde{\mathbf{y}_k} \rangle_Q = \mathbf{y}_k^\top Q \tilde{\mathbf{y}_k}
	\label{eq:inner}
\end{equation}
\begin{equation}
	\|\mathbf{u}_k\|_R = \sqrt{\mathbf{u}_k^ \top R \mathbf{u}_k}, \quad \|\mathbf{y}_k\|_Q = \sqrt{\mathbf{y}_k^ \top Q \mathbf{y}_k}
	\label{eq:norm}
\end{equation}
其中\(\tilde{\mathbf{u}_k} \in \mathbf{U}\) 和 \(\tilde{\mathbf{y}_k} \in \mathbf{Y}\) 为任意向量，$R$ 和 $Q$ 为正定加权矩阵。

基于以上预备知识，ILC致力于设计迭代学习控制律
\begin{equation}
	\mathbf{u}_{k+1}=\mathbf{u}_k+\Delta \mathbf{u}_{k+1}
	\label{kongzhiqi}
\end{equation}
实现控制输入沿迭代轴的更新，从而确保系统输出能够精确跟踪给定的期望轨迹$\mathbf{r}$，即
\begin{equation}
	\lim_{k\to\infty} \mathbf{y}_k=\mathbf{r}
	\label{mubiao}
\end{equation}
现有研究已经在线性时不变系统的迭代学习控制方面取得了丰富成果，其中基于模型的范数最优 ILC 被广泛视为性能标杆。然而这类方法高度依赖系统的精确数学模型，在实际工程中，系统模型往往难以获得：一方面，许多复杂过程（如柔性机械臂、智能材料等）机理复杂，参数难以辨识，导致建模成本高昂且精度有限；另一方面，系统参数可能随运行条件变化，使得基于固定模型的控制设计缺乏鲁棒性。因此，在模型未知或仅能获得输入输出数据的情况下，如何设计有效的迭代学习控制器成为一个重要的研究问题。为摆脱对模型的依赖，数据驱动方法应运而生。

我们首先给出一些记号，给定 \(i, j \in \mathbb{Z}\) 且 \(i \leq j\)，记 \([i, j]\) 为所有介于 \(i\) 和 \(j\) 之间的整数构成的有序集合，包含 \(i\) 和 \(j\) 本身。按惯例，若 \(i > j\)，则 \([i, j] = \emptyset\)。

设 \(i, j \in \mathbb{Z}\) 且 \(i \leq j\)。又设 \(f(i), f(i+1), \ldots, f(j) \in \mathbb{R}^n\) 是一个向量序列。定义
\begin{equation}
	f_{[i,j]} := \begin{bmatrix} f(i) & f(i+1) & \cdots & f(j-1) \end{bmatrix} \in \mathbb{R}^{n \times (j-i)}
	\label{eq:f_interval}
\end{equation}
此外，对于 \(0 \leq L \leq j-i\)，与 \(f_{[i,j]}\) 相关联的具有 \(L\) 个块行的 Hankel 矩阵定义为
\begin{equation}
	H_L(f_{[i,j]}) = 
	\begin{bmatrix}
		f(i)     & \cdots & f(j-L) \\
		\vdots   & \ddots & \vdots \\
		f(i+L-1)   & \cdots & f(j-1)
	\end{bmatrix}
	\label{eq:hankel_def}
\end{equation}
采集一段长度为 \(J\) 的输入数据 \(\mathbf{u_{d,{[0,J-1]}}} = [u_d(0), u_d(1), \dots, u_d(J-1)]^\top\) 施加于系统，并记录对应的输出数据 \(\mathbf{y_{d,{[0,J-1]}}} = [y_d(0), y_d(1), \dots, y_d(J-1)]^\top\)。
令 \( L \in [0, J-1] \)，并记由数据 \(\mathbf{u_{d,{[0,J-1]}}}, \mathbf{y_{d,{[0,J-1]}}}\) 构造的具有 \( L \) 个块行的 Hankel 矩阵为
\begin{equation}
	H_L(\mathbf{u_{d,{[0,J-1]}}}) = \begin{bmatrix}
		u_d(0)   & u_d(1)   & \cdots & u_d(J-L) \\
		u_d(1)   & u_d(2)   & \cdots & u_d(J-L+1)   \\
		\vdots   & \vdots   & \ddots & \vdots     \\
		u_d(L-1)   & u_d(L) & \cdots & u_d(J-1)
	\end{bmatrix}
	\label{eq:hankel_u}
\end{equation}
\begin{equation}
	H_L(\mathbf{y_{d,{[0,J-1]}}}) = \begin{bmatrix}
		y_d(0)   & y_d(1)   & \cdots & y_d(J-L) \\
		y_d(1)   & y_d(2)   & \cdots & y_d(J-L+1)   \\
		\vdots   & \vdots   & \ddots & \vdots     \\
		y_d(L-1)   & y_d(L) & \cdots & y_d(J-1)
	\end{bmatrix}
	\label{eq:hankel_y}
\end{equation}
若 Hankel 矩阵 \eqref{eq:hankel_u} 行满秩，则称输入数据 \( \mathbf{u_{d,{[0,J-1]}}} \) 满足阶数为 \( L \) 的持续激励条件。
将以上两个 Hankel 矩阵上下堆叠得到联合 Hankel 矩阵
\begin{equation}
	H_{L,J} := \begin{bmatrix} H_L(\mathbf{y_{d,{[0,J-1]}}}) \\ H_L(\mathbf{u_{d,{[0,J-1]}}}) \end{bmatrix}
	= \begin{bmatrix}
		y_d(0) & \cdots & y_d(J-L) \\
		y_d(1) & \cdots & y_d(J-L+1) \\
		\vdots & \ddots & \vdots \\
		y_d(L-1) & \cdots & y_d(J-1) \\
		u_d(0) & \cdots & u_d(J-L) \\
		u_d(1) & \cdots & u_d(J-L+1) \\
		\vdots & \ddots & \vdots \\
		u_d(L-1) & \cdots & u_d(J-1)
	\end{bmatrix}
	\label{eq:Hkt}
\end{equation}
我们还定义
\begin{equation}
	G_{L,J} := \begin{bmatrix} H_{L-1}(\mathbf{y_{d,{[0,J-2]}}}) \\ H_L(\mathbf{u_{d,{[0,J-1]}}}) \end{bmatrix}
	\label{eq:Gkt}
\end{equation}
注意，\( G_{L,J} \) 可以通过从 \( H_{L,J} \) 中移除最后一行的输出 \( \mathbf{y_{d,{[L-1,J-1]}}} \) 得到。
基于这些离线数据，数据驱动 ILC 问题可以描述如下：
\begin{problemnode}
	1.设计最短控制输入$\mathbf{u_{d,{[0,J-1]}}}$，确保输入输出数据 \( \mathbf{u_{d,{[0,J-1]}}} \) 和 \( \mathbf{y_{d,{[0,J-1]}}} \) 构造的联合 Hankel 矩阵的秩达到理论最大值$Lm+n$。这一秩条件保证了离线数据包含系统的全部动态信息，从而可用于后续的数据驱动控制器设计。\\
	2.给定输入输出数据 \( \mathbf{u_{d,{[0,J-1]}}} \) 和 \( \mathbf{y_{d,{[0,J-1]}}} \)。在不使用解析系统模型 \( \mathbf{G} \) 的情况下，利用输入输出数据构建形如 \eqref{kongzhiqi} 的迭代学习控制器，实现 \eqref{mubiao} 中的轨迹跟踪目标。
\end{problemnode}

\section{最短测试实验设计}

本小节针对问题 1 提出具体的数据生成方法：我们想要设计一段长度最短的输入序列 \(\mathbf{u_{d,[0,J-1]}}\)，使得由此产生的输入输出数据构造的联合 Hankel 矩阵的秩达到理论值 \(Lm + n\)。在后文的分析中，我们可以看到该条件等价于数据满足 Willems' 基本引理的要求，从而保证离线数据包含系统的全部动态信息。然而，基于 Willems' 基本引理的方法要求输入数据具有足够高的持续激励阶数，这往往导致离线数据长度过长，进而带来数据采集成本高昂、系统安全性风险增加以及后续控制计算负担沉重等问题。为此，我们可以反复利用如下引理逐步增加联合 Hankel 矩阵的秩，以实现数据长度最小化。

\begin{lemma}[The shortest experiment]
	设 \(J \geq 2\)，\(L \geq 1\)。如果
	\begin{equation}
		\operatorname{rank} G_{L,J} < m + \operatorname{rank} H_{L-1,J}
		\label{eq:cond_rank}
	\end{equation}
	则存在一个 \((m-1)\) 维仿射集 \(A_J \subseteq \mathbb{R}^m\)，使得只要 \(u_d(J) \notin A_J\)，就有
	\begin{equation}
		\operatorname{rank} H_{L,J+1} = \operatorname{rank} H_{L,J} + 1
		\label{eq:rank_increase}
	\end{equation}
	\label{lem:shortest}
\end{lemma}

\begin{proof}
	注意到 \(\ker H_{L-1,J} \times \{0_m\} \subseteq \ker G_{L,J}\)，因此 \(\dim \ker G_{L,J} \geq \dim \ker H_{L-1,J}\)。由秩-零化度定理可得
	\[
	{L}p + (L+1)m - \operatorname{rank} G_{L,J} \geq L(p+m) - \operatorname{rank} H_{L-1,J}
	\]
	即 \(\operatorname{rank} G_{L,J} \leq m + \operatorname{rank} H_{L-1,J}\)。此外，等号成立当且仅当 \(\ker H_{L-1,J} \times \{0_m\} = \ker G_{L,J}\)。因此，条件 \eqref{eq:cond_rank} 意味着存在 \(\xi_i \in \mathbb{R}^p\ (i \in [0,L-1])\) 和 \(\eta_j \in \mathbb{R}^m\ (j \in [0,L])\)，且 \(\eta_L \neq 0\)，使得
	\[
	\begin{bmatrix}
		\xi_0^\top & \xi_1^\top & \cdots & \xi_{L-1}^\top & \eta_0^\top & \eta_1^\top & \cdots & \eta_L^\top
	\end{bmatrix}
	G_{L,J} = 0
	\]
	定义集合
	\[
	\mathcal{A}_J := \left\{ v \in \mathbb{R}^m \;\middle|\; \eta_L v + \sum_{i=0}^{L-1} \xi_i^\top y(J-L+i) + \eta_i^\top u(J-L+i) = 0 \right\}
	\]
	若 \(u(J) \notin \mathcal{A}_J\)，则
	\[
	\begin{bmatrix}
		\xi_0^\top & \xi_1^\top & \cdots & \xi_{L-1}^\top & \eta_0^\top & \eta_1^\top & \cdots & \eta_L^\top
	\end{bmatrix}
	G_{L,J+1} \neq 0
	\]
	由于 \(\ker G_{L,J+1} \subseteq \ker G_{L,J}\)，由上式可得 \(\dim \ker G_{L,J+1} < \dim \ker G_{L,J}\)。因此，\(G_{L,J+1}\) 的最后一列不能表示为 \(G_{L,J}\) 的列的线性组合，进而 \(H_{L,J+1}\) 的最后一列也不能表示为 \(H_{L,J}\) 的列的线性组合，于是 \eqref{eq:rank_increase} 成立。证毕。
\end{proof}

通过该最短测试实验设计，我们可以成功解决问题 1，为后续数据驱动 ILC 提供信息量丰富且长度最短的离线数据。下一小节将基于这些数据设计迭代学习控制律，实现问题 2 的轨迹跟踪目标。

\section{数据驱动范数最优迭代学习控制}

在本节中，我们介绍 Willems' 基本引理，该引理允许将系统行为直接表示为输入输出数据的线性组合。利用这一性质，数据驱动控制方法能够实现与基于模型的控制设计相同的控制性能。本节涵盖基本概念、关键引理以及数据驱动仿真算法。需要注意的是，数据驱动控制将在行为框架下进行阐述。有关行为系统理论的更多信息，请参阅 Willems (1991) 以及 Markovsky 和 Dörfler (2021)。

首先介绍一些必要的定义。考虑系统 \eqref{eq:sys}，取定某一迭代指标$k$（为方便起见去掉下标），令长度为$L$的输入和输出超向量为
\begin{equation}
	\begin{aligned}
		\mathbf{u} &= [u(0), u(1), \dots, u(L-1)]^\top \in \mathbb{R}^{mL},\\
		\mathbf{y} &= [y(0), y(1), \dots, y(L-1)]^\top \in \mathbb{R}^{pL}
	\end{aligned}
	\label{eq:uv}
\end{equation}
长度为 \(L\) 的系统轨迹可写为
\begin{equation}
	\mathbf{w} = \begin{bmatrix} \mathbf{u} \\ \mathbf{y} \end{bmatrix} = \operatorname{col}(\mathbf{u}, \mathbf{y}) \in \mathbb{R}^{(m+p)L}
	\label{eq:trajectory}
\end{equation}
由系统 \eqref{eq:sys} 生成的所有长度为 \(L\) 的轨迹的集合构成一个子空间 \(\mathcal{G}_L\)，其定义为
\begin{equation}
	\mathcal{G}_L := \left\{ \begin{bmatrix} \mathbf{u} \\ \mathbf{y}\end{bmatrix} \in \mathbb{R}^{(m+p)L} \;\middle|\;
	\exists \left\{x(0),x(1),\cdots,x(L-1)\right\} \text{ 使得 } \eqref{eq:sys} \text{ 成立 }
	\right\}
	\label{eq:traj_space}
\end{equation}
然而需要注意，这些轨迹未必对应于同一个初始状态 $x_0$。迭代学习控制的一个关键特征在于，每次迭代开始时系统的状态会被重置为 $x_0$。在数据驱动的框架下，系统动态未知且无法直接获取状态信息。因此，我们无法显式地重置状态，而必须从输入输出数据中推断出初始条件。为此，定义系统滞后 \( lag \) 为系统的能观性指数，即使得能观性矩阵
\begin{equation}
	O = \begin{bmatrix}
		C \\
		CA \\
		\vdots \\
		CA^{lag-1}
	\end{bmatrix}
	\label{eq:obs_matrix}
\end{equation}
的秩等于 \( n \) 的最小整数 \( lag \)。系统的初始轨迹定义为
\begin{equation}
	\mathbf{w_{ini}} := \begin{bmatrix} \mathbf{u_{ini}} \\ \mathbf{y_{ini}} \end{bmatrix} \in \mathbb{R}^{(m+p)T_{ini}}
	\label{eq:initial_trajectory}
\end{equation}
其中， \( \mathbf{u_{ini}} \) 和 \( \mathbf{y_{ini}} \) 分别为初始输入和初始输出，\( T_{ini} \geq lag \) 为表示初始输入/输出响应长度的整数。此时，时刻 $T_{ini}$ 的内部状态由 $\mathbf{w_{ini}}$ 唯一确定，并可被视为后续轨迹的实际初始条件（I. Markovsky and P. Rapisarda, 2008）。\(\operatorname{col\,span}(\cdot)\) 表示矩阵的列张成的空间。

基于上述定义，我们可以定义如下输入和输出 Hankel 矩阵及其分块
\begin{equation}
	U = H_{T_{ini}+T}(\mathbf{u_{d,{[0,J-1]}}}) = \begin{bmatrix} U_p \\ U_f \end{bmatrix},
	\qquad
	Y = H_{T_{ini}+T}(\mathbf{y_{d,{[0,J-1]}}}) = \begin{bmatrix} Y_p \\ Y_f \end{bmatrix}
	\label{eq:hankel_partition}
\end{equation}
其中块 \( U_p \in \mathbb{R}^{mT_{ini} \times (J - T_{ini} - T + 1)} \), \( Y_p \in \mathbb{R}^{pT_{ini} \times (J - T_{ini} - T + 1)} \)（下标 \( p \) 表示“过去”）用于设置初始条件，而块 \( U_f \in \mathbb{R}^{mT \times (J - T_{ini} - T + 1)} \), \( Y_f \in \mathbb{R}^{pT \times (J - T_{ini} - T + 1)} \)（下标 \( f \) 表示“未来”）用于计算系统响应。下面给出数据驱动控制中的重要定理。

\begin{lemma}[Willems et al., 2005]
	考虑系统 \eqref{eq:sys} 并假设其能控。给定由系统 \eqref{eq:sys} 生成的长度为\( J \) 的数据轨迹 \( \mathbf{w_{d,{[0,J-1]}}} := \operatorname{col}(\mathbf{u_{d,{[0,J-1]}}}, \mathbf{y_{d,{[0,J-1]}}}) \)。若输入 \( \mathbf{u_{d,{[0,J-1]}}} \) 满足阶数为 \( L + n \) 的持续激励条件，则任何由系统 \eqref{eq:sys} 生成的长度为\( L \) 的轨迹 $\mathbf{w}$ 均可表示为 \( H_{L,J} \) 的列的线性组合，且任何线性组合 \( H_{L,J}g \)（其中 \( g \in \mathbb{R}^{J-L+1} \)）都是 \( \mathcal{G}_L \) 的一条轨迹，即
	\begin{equation}
		\operatorname{col\,span}(H_{L,J}) = \mathcal{G}_L
		\label{eq:willems}
	\end{equation}
	\label{lem:willems}
\end{lemma}

引理 \ref{lem:willems} 表明，系统行为可由满足持续激励条件的已有数据来表示，其有一个重要应用是数据驱动仿真，如下面的算法所示（Markovsky and Rapisarda, 2008）。

\begin{algorithm}[H]
	\SetAlgoLined
	\SetAlgorithmName{算法}{算法}{算法目录}
	\SetKwInput{KwIn}{输入}
	\SetKwInput{KwOut}{输出}
	\SetKw{Return}{返回}
	\caption{数据驱动仿真}
	\label{alg:data_driven_simulation}
	\KwIn{数据轨迹 $\mathbf{w_{d,{[0,J-1]}}}$，初始轨迹 $\mathbf{w_{ini}}$，输入 $\mathbf{u}$}
	\KwOut{系统输出 $\mathbf{y}$}
	求解矩阵方程：
	\begin{equation}
		\begin{bmatrix} U_p \\ Y_p \\ U_f \end{bmatrix} g = \begin{bmatrix} \mathbf{u_{ini}} \\ \mathbf{y_{ini}} \\ \mathbf{u} \end{bmatrix}
		\label{eq:sim_g}
	\end{equation}
	计算输出：
	\begin{equation}
		\mathbf{y} = Y_f g
		\label{eq:sim_y}
	\end{equation}
	\Return $\mathbf{y}$
\end{algorithm}
需要注意的是，所有具有零初始条件的轨迹 \( \mathbf{w} \in \mathcal{G}_L \) 的集合可以构成系统 \eqref{eq:sys} 的一个子行为，记为 \( \mathcal{G}_L^0 \)。该子行为的一组基可以由矩阵 \( W_0 \)的列给出，使得
\begin{equation}
	\operatorname{col\,span}(W_0) = \mathcal{G}_L^0
	\label{eq:zero_initial_subspace}
\end{equation}
其中 \( W_0 \) 可通过以下算法计算。

\begin{algorithm}[H]
	\SetAlgoLined
	\SetAlgorithmName{算法}{算法}{算法目录}
	\SetKwInput{KwIn}{输入}
	\SetKwInput{KwOut}{输出}
	\SetKw{Return}{返回}
	\caption{零初始条件子行为的数据驱动矩阵计算}
	\label{alg:zero_initial_matrix}
	\KwIn{数据轨迹 $\mathbf{w_{d,{[0,J-1]}}}$}
	\KwOut{零初始条件子行为 $\mathcal{G}_L^0$ 的数据驱动矩阵表示 $W_0$}
	求解以下方程：
	\begin{equation}
		\begin{bmatrix} U_p \\ Y_p \\ U_f \end{bmatrix} g =
		\begin{bmatrix}
			0_{mT_{\mathrm{ini}} \times (J - T + 1)} \\
			0_{pT_{\mathrm{ini}} \times (J - T + 1)} \\
			H_L(\mathbf{u_{d,{[0,J-1]}}})
		\end{bmatrix}
		\label{eq:solve_g}
	\end{equation}
	其中 $0_{m,n}$ 表示 $m \times n$ 的零矩阵。所得结果用于计算
	\begin{equation}
		Y_f g = Y_0
		\label{eq:Y0}
	\end{equation}
	于是，$W_0$ 可写为
	\begin{equation}
		W_0 = \begin{bmatrix} H_L(\mathbf{u_{d,{[0,J-1]}}}) \\ Y_0 \end{bmatrix}
		\label{eq:W_0}
	\end{equation}
	\Return $W_0$
\end{algorithm}
基于模型的范数最优ILC算法中的优化问题利用模型信息来表达系统的输入输出关系，即 \( y = Gu + d \)。如前所述，在实际中辨识这一解析系统模型 \( G \) 可能很困难。因此，在数据驱动范数最优ILC算法中，应用 Willems' 基本引理，利用已有数据来表达系统的输入输出关系，从而无需解析模型，具体算法如下所示。

\begin{algorithm}[H]
	\SetAlgoLined
	\SetAlgorithmName{算法}{算法}{算法目录}
	\SetKwInput{KwIn}{输入}
	\SetKwInput{KwOut}{输出}
	\SetKw{Return}{返回}
	\caption{数据驱动范数最优迭代学习控制}
	\label{alg:data_driven_NOILC}
	\KwIn{数据轨迹 $\mathbf{w_{d,{[0,J-1]}}} = \operatorname{col}(\mathbf{u_{d,{[0,J-1]}}}, \mathbf{y_{d,{[0,J-1]}}})$，其中 $\mathbf{u_{d,{[0,J-1]}}}$ 满足阶数为 $T+T_{\mathrm{ini}}+n$ 的持续激励条件；参考轨迹 $\mathbf{r}$；权重矩阵 $Q$ 和 $R$；上一次的输入 $\mathbf{u_k}$ 和误差 $\mathbf{e_k}$}
	\KwOut{下一次的输入 $\mathbf{u_{k+1}}$}
	求解如下优化问题：
	\begin{equation}、
		\mathbf{u_{k+1}} = \arg \min_{\mathbf{u_{k+1}}} \|\mathbf{e_{k+1}}\|_Q^2 + \|\mathbf{u_{k+1}} - \mathbf{u_k}\|_R^2
		\label{eq:opt_noilc}
	\end{equation}
	满足约束
	\begin{equation}
		\begin{bmatrix} U_p \\ Y_p \\ U_f \\ Y_f \end{bmatrix} (g_{k+1}-g_k) = \begin{bmatrix} 0\\ 0 \\ \mathbf{u_{k+1}-u_k} \\ \mathbf{y_{k+1}-y_k} \end{bmatrix}, \quad \mathbf{e_{k+1}} = \mathbf{r} - \mathbf{y_{k+1}}
		\label{eq:constraint_noilc}
	\end{equation}
	\Return $\mathbf{u_{k+1}}$
\end{algorithm}
\begin{proposition}
	算法 \ref{alg:data_driven_NOILC} 与基于模型的范数最优ILC算法等价。
	\label{prop:noilc_equivalence}
\end{proposition}

\begin{proof}
	为证明该命题，我们展示数据驱动范数最优ILC算法的优化问题等价于基于模型的范数最优ILC优化问题 \eqref{eq:model_based_NOILC}。显然，\eqref{eq:opt_noilc} 的代价函数与基于模型的范数最优ILC优化问题 \eqref{eq:cost_NOILC} 相同。因此，我们将证明 \eqref{eq:constraint_noilc} 的约束部分也与基于模型的范数最优ILC优化问题一致。
	
	定义虚拟初始输入轨迹 \( \mathbf{u_{ini}} \) 和虚拟初始输出轨迹 \( \mathbf{y_{ini}} \)，使得它们满足以下方程
	\begin{equation}
		\begin{bmatrix} U_p \\ Y_p \\ U_f \\ Y_f \end{bmatrix} g = \begin{bmatrix} \mathbf{u_{ini}} \\ \mathbf{y_{ini}} \\ 0 \\ \mathbf{d} \end{bmatrix}
		\label{eq:virtual_initial}
	\end{equation}
	上述方程表明，当系统的初始输入输出轨迹为 \( \mathbf{u_{ini}} \) 和 \( \mathbf{y_{ini}} \) 时，系统的自由响应为 \( \mathbf{d} \)。需要注意的是，定义虚拟初始轨迹的原因是：在ILC系统中，当时间到达 \( N+1 \) 时，系统状态直接被重置为 \( x_0 \)，因此ILC系统可能没有实际的初始轨迹。于是存在 \( g_{k+1} \) 和 \( g_k \) 使得以下方程成立
	\begin{equation}
		\begin{bmatrix} U_p \\ Y_p \\ U_f \\ Y_f \end{bmatrix} g_{k+1} = \begin{bmatrix} \mathbf{u_{ini}} \\ \mathbf{y_{ini}} \\ \mathbf{u_{k+1}} \\ \mathbf{y_{k+1}} \end{bmatrix}, \quad
		\begin{bmatrix} U_p \\ Y_p \\ U_f \\ Y_f \end{bmatrix} g_{k} = \begin{bmatrix} \mathbf{u_{ini}} \\ \mathbf{y_{ini}} \\ \mathbf{u_k} \\ \mathbf{y_k} \end{bmatrix}
		\label{eq:g1_g2}
	\end{equation}
	将两式相减，可得
	\begin{equation}
		\begin{bmatrix} U_p \\ Y_p \\ U_f \\ Y_f \end{bmatrix} (g_{k+1}-g_k) = \begin{bmatrix} 0 \\ 0 \\ \mathbf{u_{k+1}} - \mathbf{u_k} \\ \mathbf{y_{k+1}} - \mathbf{y_k} \end{bmatrix}
		\label{eq:g3}
	\end{equation}
	利用 Willems' 基本引理，上述方程等价于基于模型的输入输出关系
	\begin{equation}
		\mathbf{y_{k+1}} - \mathbf{y_k} = \mathbf{G} (\mathbf{u_{k+1}} - \mathbf{u_k})
		\label{eq:diff_io}
	\end{equation}
	将 \( \mathbf{y_k} = \mathbf{G} \mathbf{u_k} + \mathbf{d} \) 代入上式，可得
	\begin{equation}
		\mathbf{y_{k+1}} = \mathbf{G} \mathbf{u_{k+1}} + \mathbf{d}
		\label{eq:io_relation}
	\end{equation}
	由此可知，\eqref{eq:constraint_noilc} 的约束部分等价于基于模型的范数最优ILC算法的约束部分。
\end{proof}

\begin{proposition}
	\eqref{eq:opt_noilc} \eqref{eq:constraint_noilc} 描述的优化问题的解析解由下式给出：
	\begin{equation}
		\mathbf{u_{k+1}} = \mathbf{u_k} + \begin{bmatrix} I_{mL} \\ 0_{pL,mL} \end{bmatrix}^\top W_0 (W_0^\top S W_0)^+ W_0^\top S \begin{bmatrix} 0_{mL,1} \\ \mathbf{e_k} \end{bmatrix},
		\label{eq:solution_noilc}
	\end{equation}
	其中 \( I_m \) 表示 \( m \times m \) 单位矩阵；\( 0_{m,n} \) 表示 \( m \times n \) 零矩阵；\( + \) 表示伪逆；\( S = \begin{bmatrix} R & 0 \\ 0 & Q \end{bmatrix} \)；\( W_0 \) 为零初始条件响应矩阵，可通过算法 \ref{alg:zero_initial_matrix} 求得（Markovsky and Rapisarda, 2008）。
\end{proposition}

\begin{proof}
	零初始条件响应矩阵 \( W_0 \) 具有以下性质（Markovsky and Rapisarda, 2008）：
	\begin{equation}
		\operatorname{col\,span}(W_0) = \operatorname{col\,span}\left( \begin{bmatrix} I \\ \mathbf{G} \end{bmatrix} \right)
		\label{eq:W0_property}
	\end{equation}
	因此，对于任意输入差值 \( \mathbf{u_{k+1}} - \mathbf{u_k} \)，存在 \( g_k \) 使得
	\begin{equation}
		W_0 g_k = \begin{bmatrix} I \\ \mathbf{G} \end{bmatrix} (\mathbf{u_{k+1}} - \mathbf{u_k}) = \begin{bmatrix} \mathbf{u_{k+1}} - \mathbf{u_k} \\ \mathbf{y_{k+1}} - \mathbf{y_k} \end{bmatrix}
		\label{eq:W0_g}
	\end{equation}
	利用上式，优化问题 \eqref{eq:opt_noilc} 和 \eqref{eq:constraint_noilc} 可重写为
	\begin{equation}
		\mathbf{u_{k+1}} = \arg \min_{\mathbf{u_{k+1}}} \| \mathbf{e_{k+1}} \|_Q^2 + \| \mathbf{u_{k+1}} - \mathbf{u_k} \|_R^2
	\end{equation}
	满足约束
	\begin{equation}
		\begin{bmatrix} \mathbf{u_{k+1}} - \mathbf{u_k} \\ \mathbf{y_{k+1}} - \mathbf{y_k} \end{bmatrix} = W_0 g_k, \quad \mathbf{e_{k+1}} = \mathbf{r} - \mathbf{y_{k+1}}
		\label{eq:constrained_opt}
	\end{equation}
	注意到 \( \mathbf{e_{k+1}} = \mathbf{e_k} - (\mathbf{y_{k+1}} - \mathbf{y_k}) \)，且 \( \mathbf{e_k} = \mathbf{r} - \mathbf{y_k} \)。令 \( \mathbf{\Delta u} = \mathbf{u_{k+1}} - \mathbf{u_k} \)，\( \mathbf{\Delta y} = \mathbf{y_{k+1}} - \mathbf{y_k} \)。由 \eqref{eq:W0_g} 知 \( \mathbf{\Delta u} = [I \ 0] W_0 g_k \)，\( \mathbf{\Delta y} = [0 \ I] W_0 g_k \)。则代价函数可表示为关于 \( g_k \) 的二次型：
	\begin{align*}
		J(g_k) &= \| \mathbf{e_k} - \mathbf{\Delta y} \|_Q^2 + \| \mathbf{\Delta u} \|_R^2 \\
		&= \left( \mathbf{e_k} - [0 \ I] W_0 g_k \right)^\top Q \left( \mathbf{e_k} - [0 \ I] W_0 g_k \right) + \left( [I \ 0] W_0 g_k \right)^\top R \left( [I \ 0] W_0 g_k \right)
	\end{align*}
	展开并忽略常数项，得
	\[
	J(g_k) = -2 \mathbf{e_k}^\top Q [0 \ I] W_0 g_k + g_k^\top W_0^\top \begin{bmatrix} R & 0 \\ 0 & Q \end{bmatrix} W_0 g_k + Const
	\]
	记 \( S = \begin{bmatrix} R & 0 \\ 0 & Q \end{bmatrix} \)，则 \( J(g_k) = g_k^\top (W_0^\top S W_0) g_k - 2 \mathbf{e_k}^\top [0 \ I] W_0 g_k + Const \)。对 \( g_k \) 求导并令导数为零得：
	\[
	2 W_0^\top S W_0 g_k - 2 W_0^\top \begin{bmatrix} 0 \\ I \end{bmatrix} \mathbf{e_k} = 0,
	\]
	即
	\[
	W_0^\top S W_0 g_k = W_0^\top \begin{bmatrix} 0 \\ \mathbf{e_k} \end{bmatrix}
	\]
	由于 \( W_0^\top S W_0 \) 可能奇异，其最小二乘解为
	\[
	g_k = (W_0^\top S W_0)^+ W_0^\top \begin{bmatrix} 0 \\ \mathbf{e_k} \end{bmatrix}
	\]
	代入 \( \mathbf{\Delta u} = [I \ 0] W_0 g_k \)，得
	\begin{equation}
		\mathbf{\Delta u} = [I \ 0] W_0 (W_0^\top S W_0)^+ W_0^\top \begin{bmatrix} 0 \\ \mathbf{e_k} \end{bmatrix}
		\label{eq:delta_u}
	\end{equation}
	因此，
	\begin{equation}
		\mathbf{u_{k+1}} = \mathbf{u_k} + \begin{bmatrix} I \\ 0 \end{bmatrix}^\top W_0 (W_0^\top S W_0)^+ W_0^\top S \begin{bmatrix} 0 \\ \mathbf{e_k} \end{bmatrix}
		\label{eq:solution_step}
	\end{equation}
	证毕。
\end{proof}